\documentclass[11pt]{article}

\usepackage{amsmath}
\usepackage{amssymb}
\usepackage{amsthm}
\usepackage{geometry}
\usepackage{hyperref}
\geometry{margin=1in}

\newtheorem{theorem}{Theorem}
\newtheorem{remark}{Remark}

\title{Braid-Group Diagnostics and Non-Commutative Execution Ordering}
\author{Utah Finance Library --- Continuous-Time Topological Allocation Suite}
\date{}

\begin{document}
\maketitle

\noindent
This note documents \texttt{src/core/braid\_router.py}: a Temperley--Lieb /
Kauffman-bracket evaluator for the topological complexity of an execution
schedule, and an optimiser for executing a basket under \emph{non-commutative}
market impact.

\section*{Theorem (knot invariant of an execution braid)}

\begin{theorem}[Kauffman bracket via the Temperley--Lieb algebra]
Let a sequence of pairwise reorderings of $n$ simultaneous child orders be a word
in the Artin braid group $B_n$. Map each generator into the Temperley--Lieb
algebra $TL_n(\delta)$, $\delta = -A^2 - A^{-2}$, by the Kauffman skein
\begin{equation}
\sigma_i \;\mapsto\; A\,\mathbf{1} + A^{-1} e_i,
\qquad
\sigma_i^{-1} \;\mapsto\; A^{-1}\,\mathbf{1} + A\,e_i,
\end{equation}
where $e_i$ are the TL generators satisfying $e_i^2 = \delta e_i$,
$e_i e_{i\pm1} e_i = e_i$, $e_i e_j = e_j e_i$ for $|i-j|\ge 2$. Then the Markov
trace of the resulting element equals the Kauffman bracket of the braid closure,
and the writhe-normalised quantity $f_L = (-A^3)^{-w}\langle L\rangle$ is the
Jones polynomial $V_L(t)$ at $t = A^{-4}$, an isotopy invariant of the closure.
\end{theorem}

\begin{proof}
The skein relations make $\langle\cdot\rangle$ invariant under Reidemeister~II
and~III; the loop value $\delta$ implements removal of disjoint circles. The
$TL_n$ relations are exactly the diagram-composition rules (planar matchings,
with each closed loop contributing a factor $\delta$), so the homomorphism
$B_n \to TL_n$ is well defined. Closing the braid and counting loops realises the
Markov trace, giving $\langle L\rangle$. Multiplying by $(-A^3)^{-w}$ cancels the
Reidemeister~I anomaly, yielding the (oriented) isotopy invariant $V_L(A^{-4})$.
$\square$
\end{proof}

\section*{Execution ordering under non-commutative impact}

Let $L\in\mathbb{R}^{m\times m}_{\ge 0}$ be a cross-impact matrix where $L_{ab}$ is
the extra slippage incurred on trade $b$ if $a$ executes before it. For an ordering
$\pi$ the total modelled impact is
\begin{equation}
C(\pi) = \sum_{a \text{ before } b \text{ in } \pi} L_{ab}.
\end{equation}
Because $L$ is asymmetric, $C$ depends on order ($a$-then-$b$ differs from
$b$-then-$a$): this is the linear-ordering problem, solved exactly by enumeration
for small baskets and greedily otherwise. The realised reordering is realised
concretely as a braid word whose free-reduced length is its crossing number.

\section*{Implementation correspondence}

\begin{itemize}
\item \texttt{free\_reduce}, \texttt{crossing\_number}, \texttt{braid\_permutation},
\texttt{writhe}: braid-word combinatorics.
\item \texttt{kauffman\_bracket}, \texttt{jones\_polynomial\_value}: the $TL_n$
diagram algebra and Markov trace (the tests verify $e_i^2=\delta e_i$,
$e_1 e_2 e_1 = e_1$, the unknot normalisation, Reidemeister~II, and the trefoil
Jones polynomial).
\item \texttt{optimize\_execution\_braid}, \texttt{ordering\_impact\_cost}:
minimisation of $C(\pi)$, with the saving versus the submitted order routed
through the transparent universal tithe.
\end{itemize}

\section*{Scope of the claim (honest framing)}

\begin{remark}
The Jones/Kauffman value is an exact topological \emph{complexity descriptor} of
an execution schedule, not a P\&L. The ordering optimiser minimises a
\emph{modelled} cost over permutations.
\end{remark}

\begin{remark}
There is \textbf{no} ``zero-slippage'' guarantee and no reduction of impact to an
``absolute infimum''. The output is only as good as the cross-impact matrix $L$,
which must be calibrated and validated on your own fills; a mis-specified $L$
yields a mis-ordered basket. The optimiser reduces modelled cost relative to the
naive order --- it does not make trading frictionless.
\end{remark}

\end{document}
